\documentclass[aspectratio=169,11pt]{beamer}
% Shared Beamer style for practice contest solution walkthroughs.
\usepackage[utf8]{inputenc}
\usepackage[T1]{fontenc}
\usepackage{lmodern}
\usepackage{amsmath,amssymb}
\usepackage{xcolor}
\usepackage{hyperref}
\usepackage{listings}

\definecolor{CPNavy}{HTML}{172033}
\definecolor{CPBlue}{HTML}{2563EB}
\definecolor{CPGreen}{HTML}{16A34A}
\definecolor{CPOrange}{HTML}{F97316}
\definecolor{CPGray}{HTML}{64748B}
\definecolor{CPLight}{HTML}{F8FAFC}
\definecolor{CPCodeBg}{HTML}{F1F5F9}
\definecolor{CPCodeRule}{HTML}{CBD5E1}

\usetheme{default}
\usefonttheme{professionalfonts}
\setbeamertemplate{navigation symbols}{}
\setbeamersize{text margin left=0.65cm,text margin right=0.65cm}
\setbeamertemplate{itemize item}{\color{CPBlue}$\blacktriangleright$}
\setbeamertemplate{itemize subitem}{\color{CPGray}$\triangleright$}

\setbeamercolor{normal text}{fg=CPNavy,bg=white}
\setbeamercolor{frametitle}{fg=white,bg=CPNavy}
\setbeamercolor{title}{fg=white,bg=CPNavy}
\setbeamercolor{subtitle}{fg=CPLight,bg=CPNavy}
\hypersetup{colorlinks=true,urlcolor=CPBlue}

\setbeamertemplate{frametitle}{%
  \nointerlineskip%
  \begin{beamercolorbox}[wd=\paperwidth,ht=0.88cm,dp=0.22cm,leftskip=0.55cm,rightskip=0.35cm]{frametitle}
    \usebeamerfont{frametitle}\insertframetitle\hfill\small\insertframenumber/\inserttotalframenumber
  \end{beamercolorbox}%
}

\setbeamertemplate{footline}{%
  \leavevmode%
  \hbox{%
    \begin{beamercolorbox}[wd=.58\paperwidth,ht=2.4ex,dp=1ex,leftskip=.5cm]{author in head/foot}%
      \color{CPGray}\insertshorttitle
    \end{beamercolorbox}%
    \begin{beamercolorbox}[wd=.42\paperwidth,ht=2.4ex,dp=1ex,rightskip=.5cm plus1fil]{date in head/foot}%
      \hfill\color{CPGray}\insertshortauthor
    \end{beamercolorbox}%
  }%
}

\setbeamertemplate{title page}{%
  \vspace*{1.25cm}
  \begin{beamercolorbox}[wd=\paperwidth,sep=0.65cm,leftskip=0.15cm,rightskip=0.15cm]{title}
    {\color{white}\Huge\bfseries\inserttitle\par}
    \vspace{0.35cm}
    {\color{CPLight}\Large\insertsubtitle\par}
    \vspace{0.6cm}
    {\color{CPLight}\large\insertauthor\par}
  \end{beamercolorbox}
  \vfill
}

\newcommand{\sectiontag}[1]{\textbf{\color{CPBlue}#1}}
\newenvironment{tightitemize}{\begin{itemize}\small\setlength{\itemsep}{0.18cm}}{\end{itemize}}
\lstdefinestyle{cpcode}{
  language=Python,
  basicstyle=\ttfamily\tiny,
  keywordstyle=\color{CPBlue}\bfseries,
  commentstyle=\color{CPGray}\itshape,
  stringstyle=\color{CPGreen},
  backgroundcolor=\color{CPCodeBg},
  frame=single,
  framerule=0.25pt,
  rulecolor=\color{CPCodeRule},
  showstringspaces=false,
  columns=fullflexible,
  keepspaces=true,
  breaklines=true,
  breakatwhitespace=false,
  tabsize=4,
  xleftmargin=0.08cm,
  xrightmargin=0.08cm,
  aboveskip=0.08cm,
  belowskip=0.08cm
}


\title[beepers]{Collecting Beepers}
\subtitle{Day 3 solution walkthrough}
\author[Solution Deck]{Competitive Programming Bootcamp}
\date{}

\begin{document}

\begin{frame}[plain]
  \titlepage
\end{frame}

% BEGIN GENERATED STATEMENT INTERPRETATION
\begin{frame}{Interpret the Word Problem}
\small
\sectiontag{Statement excerpts:}
\begin{tightitemize}
  \item \textit{``each of the beepers, and return back again''}
  \item \textit{``never be more than \$10\$ beepers''}
\end{tightitemize}

\vspace{0.18cm}
\sectiontag{Programming interpretation:}
\begin{tightitemize}
  \item The grid size is irrelevant once coordinates are known.
  \item With at most 10 beepers, bitmask dynamic programming over visited sets is feasible.
  \item The state is current point plus the set of collected beepers.
\end{tightitemize}
\end{frame}
% END GENERATED STATEMENT INTERPRETATION







\begin{frame}{Problem Model}
\small
\sectiontag{Textbook connection:} CP4 3.5 e. Travelling Salesman Problem

\vspace{0.25cm}
\sectiontag{Core technique:} Bitmask DP for small TSP

\vspace{0.25cm}
\begin{tightitemize}
  \item Start at the robot's home position.
  \item Visit every beeper exactly once in any order.
  \item Return to the starting position with minimum Manhattan distance.
\end{tightitemize}
\end{frame}

\begin{frame}{How to Recognize the Approach}
\begin{tightitemize}
  \item The number of beepers is small, but all visit orders are factorial.
  \item A bitmask compactly records which beepers are already collected.
  \item The current location plus mask fully determines the remaining problem.
\end{tightitemize}
\end{frame}

\begin{frame}{Algorithm}
\begin{tightitemize}
  \item Store points[0] as the start and points[1..B] as beepers.
  \item Define dp(current, mask) as minimum remaining route length.
  \item If mask has all B bits set, return distance back to the start.
  \item Otherwise try every unvisited beeper, pay travel distance, and recurse.
  \item Memoize dp with lru\_cache.
\end{tightitemize}
\end{frame}

\begin{frame}{Walkthrough of the Key State}
\begin{tightitemize}
  \item Bit i of mask corresponds to beeper i at points[i + 1].
  \item new\_mask = mask | (1 << i) marks a beeper as visited.
  \item Manhattan distance is abs(dx) + abs(dy).
\end{tightitemize}
\end{frame}

\begin{frame}{Why the Algorithm Is Correct}
\begin{tightitemize}
  \item Every feasible tour has a first unvisited beeper chosen from the current state.
  \item The recurrence tries every such next beeper and takes the minimum.
  \item Memoization does not change the recurrence; it only avoids recomputation.
\end{tightitemize}
\end{frame}

\begin{frame}{Complexity and Implementation Notes}
\small
\sectiontag{Complexity:} O(B\textasciicircum{}2 * 2\textasciicircum{}B) time and O(B * 2\textasciicircum{}B) memory.

\vspace{0.3cm}
\begin{tightitemize}
  \item The grid dimensions are read but not needed for distance calculations.
  \item This is the classic small TSP state: position plus visited set.
\end{tightitemize}
\end{frame}



% BEGIN GENERATED CODE WALKTHROUGH
\begin{frame}[fragile]{Code: Distance Helper}
\scriptsize
\sectiontag{Source lines:} \texttt{beepers.py:45--52}

\vspace{0.08cm}
\begin{columns}[T,totalwidth=\textwidth]
\begin{column}{0.58\textwidth}
\begin{lstlisting}[style=cpcode]
import sys
from functools import lru_cache

def distance(a, b):
    return abs(a[0] - b[0]) + abs(a[1] - b[1])
\end{lstlisting}
\end{column}
\begin{column}{0.38\textwidth}
\sectiontag{What this chunk does}
\begin{tightitemize}
  \item Karel moves only horizontally and vertically.
  \item Manhattan distance is therefore the exact travel cost between two coordinates.
\end{tightitemize}
\end{column}
\end{columns}
\end{frame}

\begin{frame}[fragile]{Code: Read Points and Goal Mask}
\scriptsize
\sectiontag{Source lines:} \texttt{beepers.py:54--76}

\vspace{0.08cm}
\begin{columns}[T,totalwidth=\textwidth]
\begin{column}{0.58\textwidth}
\begin{lstlisting}[style=cpcode]
def solve_case():
    width, height = map(int, input().split())

    start_x, start_y = map(int, input().split())

    num_beepers = int(input())

    points = [(start_x, start_y)]

    for _ in range(num_beepers):
        x, y = map(int, input().split())
        points.append((x, y))

    all_visited = (1 << num_beepers) - 1
\end{lstlisting}
\end{column}
\begin{column}{0.38\textwidth}
\sectiontag{What this chunk does}
\begin{tightitemize}
  \item points[0] is the starting location; later entries are beepers.
  \item The all\_visited mask has one bit set for each beeper.
  \item When the current mask equals all\_visited, collection is complete.
\end{tightitemize}
\end{column}
\end{columns}
\end{frame}

\begin{frame}[fragile]{Code: Memoized DP Base Case}
\scriptsize
\sectiontag{Source lines:} \texttt{beepers.py:78--95}

\vspace{0.08cm}
\begin{columns}[T,totalwidth=\textwidth]
\begin{column}{0.58\textwidth}
\begin{lstlisting}[style=cpcode]
@lru_cache(None)
def dp(current, mask):

    if mask == all_visited:
        return distance(points[current], points[0])
\end{lstlisting}
\end{column}
\begin{column}{0.38\textwidth}
\sectiontag{What this chunk does}
\begin{tightitemize}
  \item lru\_cache stores answers for repeated current/mask states.
  \item The base case returns directly to the starting point.
  \item This turns the factorial route search into a bitmask DP.
\end{tightitemize}
\end{column}
\end{columns}
\end{frame}

\begin{frame}[fragile]{Code: Try Unvisited Beepers}
\scriptsize
\sectiontag{Source lines:} \texttt{beepers.py:97--124}

\vspace{0.08cm}
\begin{columns}[T,totalwidth=\textwidth]
\begin{column}{0.58\textwidth}
\begin{lstlisting}[style=cpcode]
    best = float("inf")

    for beeper in range(num_beepers):

        if mask & (1 << beeper):
            continue

        next_point = beeper + 1

        new_mask = mask | (1 << beeper)

        travel = distance(points[current], points[next_point])

        total = travel + dp(next_point, new_mask)

        best = min(best, total)

    return best

return dp(0, 0)
\end{lstlisting}
\end{column}
\begin{column}{0.38\textwidth}
\sectiontag{What this chunk does}
\begin{tightitemize}
  \item The loop skips any beeper whose bit is already set.
  \item new\_mask marks the chosen beeper as collected.
  \item The best recursive choice is the minimum travel plus future cost.
\end{tightitemize}
\end{column}
\end{columns}
\end{frame}

\begin{frame}[fragile]{Code: Run Each Scenario}
\scriptsize
\sectiontag{Source lines:} \texttt{beepers.py:127--138}

\vspace{0.08cm}
\begin{columns}[T,totalwidth=\textwidth]
\begin{column}{0.58\textwidth}
\begin{lstlisting}[style=cpcode]
def main():
    global input
    input = sys.stdin.buffer.readline

    test_cases = int(input())

    for _ in range(test_cases):
        answer = solve_case()
        print(answer)

main()
\end{lstlisting}
\end{column}
\begin{column}{0.38\textwidth}
\sectiontag{What this chunk does}
\begin{tightitemize}
  \item The main loop processes every Karel scenario independently.
  \item solve\_case returns the minimum round-trip distance.
\end{tightitemize}
\end{column}
\end{columns}
\end{frame}
% END GENERATED CODE WALKTHROUGH

\begin{frame}{Source}
\small
\sectiontag{Solution file:} \texttt{practice\_contests/day\_3/beepers.py}

\vspace{0.3cm}
\sectiontag{Problem:} \url{https://open.kattis.com/problems/beepers}

\vspace{0.45cm}
Use this deck with the source code open beside it. The slides explain the reasoning; the code shows the exact input handling and output formatting.
\end{frame}

\end{document}
